\documentclass[12pt]{article}
\usepackage[a4paper,margin=1in]{geometry}
\usepackage[utf8]{inputenc}
\usepackage[T1]{fontenc}
\usepackage{amsmath,amssymb}
\usepackage{textcomp}
\usepackage{parskip}
\begin{document}
\begin{center}
{\Large \textbf{2007 STEP III}}
\end{center}

\subsection*{Page 1}
\begin{flushleft}
\bigskip\textbf{Section A: Pure Mathematics}\par
\bigskip\textbf{Question 1.} In this question, do not consider the special cases in which the denominators of any of your\par
expressions are zero.\par
\[Express tan ( 01 + 02+ 03 + 04) in terms of ti, where t^1 tan01, etc.\]
\[=\]
Given that tan 01, tan 02, tan 03 and tan 04 are the four roots of the equation\par
\[at + bt3 + ct² + dt +e\]
\[=\]
(where a ≠0), find an expression in terms of a, b, c, d and e for tan(01+ 02 + 03+ 04).\par
The four real numbers 01, 02, 03 and 04 lie in the range 0 < A; < 2\pi and satisfy the equation\par
\[pcos 20+ cos(0- a) +p= 0,\]
\[where p and a are independent of 0. Show that 01+02+ 03 +01 = nn for some integer n.\]
(2n)!\par
\bigskip\textbf{Question 2.} (i) Show that 1.3.5.7. ... .(2n\par
-\par
1)\par
\[=\]
\[2n! and that, for <\]
1 (2n)!\par
\[=1+2 (n!)2\]
\[\sqrt\{1\}-4x\]
\[n=1\]
\textbf{(ii)} By differentiating the above result, deduce that\par
(2n)!\par
Σ\par
60.\par
\[n! (n 1)! 25 =\]
\[n=1 -\]
\textbf{(iii)} Show that\par
2n+1(2n)!\par
1.\par
\[32n (n+1)! n! =\]
\[OCR 2007 9475 Jun07\]
\end{flushleft}
\subsection*{Page 2}
\begin{flushleft}
\bigskip\textbf{Question 3.} A sequence of numbers, Fi, F2,..., is defined by F₁ 1, F2 1, and\par
\[= =\]
\[Fn= Fn-1+ Fn-2 for n \geq 3.\]
\textbf{(i)} Write down the values of F3, FA,., Fs.\par
\textbf{(ii)} Prove that F2k43 F2k+1- F2k+2 -F2*+2F2k+ Fk+1\par
\[=\]
\textbf{(iii)} Prove by induction or otherwise that Fan+1F2n-1- F = 1 and deduce that Fon +1\par
\[is divisible by F2n+1\]
\textbf{(iv)} Prove that Fn-1+1 is divisible by F2n+1\par
\bigskip\textbf{Question 4.} A curve is given parametrically by\par
a(cost+In t),\par
\[x =\]
tan\par
\[y = a sint,\]
where 0 < t < and a is a positive constant. Show that d d x y = = taut and sketch the curve.\par
Let P be the point with parameter t and let Q be the point where the tangent to the curve\par
\[at P meets the x-axis. Show that PQ = a.\]
The radius of curvature, p, at P is defined by\par
(22)\par
\[p=\]
j- yä\par
\[where the dots denote differentiation with respect to t. Show that p = a cot t.\]
The point C lies on the normal to the curve at P, a distance p from P and above the curve.\par
Show that CQ is parallel to the y-axis.\par
\bigskip\textbf{Question 5.} Let y = In(x2-1), where x > 1, and let r and 0 be functions of x determined by r = \sqrtx² - 1\par
\[and coth0 = x. Show that\]
dy 2 cosh 0 d²y 2 cosh 20\par
da\par
and\par
\[dx2\]
\[==\]
\[p^2\]
d³y\par
\[and find an expression in terms of r and 0 for d^23\]
dny\par
Find, with proof, a similar formula ford d an in terms of r and 0.\par
[Turn\par
\[OCR 2007 9475 Jun07 over\]
\end{flushleft}
\subsection*{Page 3}
\begin{flushleft}
\bigskip\textbf{Question 6.} The distinct points P, Q, R and S in the Argand diagram lie on a circle of radius a centred at\par
the origin and are represented by the complex numbers p, q, r and s, respectively. Show that\par
\[=-a2P-9\]
pq\par
p* q*\par
-\par
\[Deduce that, if the chords PQ and RS are perpendicular, then pq + rs = 0.\]
The distinct points A1, A2,..., An (where n > 3) lie on a circle. The points B1, B2, .., Bn lie\par
on the same circle and are chosen so that the chords B1B2, B2B3,..., B B1 are perpendicular,\par
\[respectively, to the chords A^1 A^2, A2A3,..., AnA. Show that, for n = 3, there are only two\]
choices of B₁ for which this is possible. What is the corresponding result for n = 4? State the\par
corresponding results for values of n greater than 4.\par
\bigskip\textbf{Question 7.} The functions s(x) (0<x < 1) and t(x) (x> 0), and the real number p, are defined by\par
1 1 1\par
\[(2), (2) = Jo de =2 1+42 du.\]
For this question, do not evaluate any of the above integrals explicitly in terms of inverse\par
trigonometric functions or the number .\par
100 1\par
fe\par
\textbf{(i)} Use the substitution u = u to show that t(r) = J1/x 1 1+ v2 du. Hence evaluate\par
\[t(1/1) + t(x) in terms of p and deduce that 2t(1) = p.\]
น du 1\par
\textbf{(ii)} Let y= \sqrt1+42 Express u in terms of y, and show that dy = = \sqrt(1-у2)3\par
By making a substitution in the integral for t(x), show that\par
(2)(7)\par
Deduce that s(2) p\par
\[=\]
\textbf{(iii)} Let z = 1-\par
+\par
\[\sqrt\{3\} 3 u Show that t() 13\]
\[\sqrt\{3\}\]
\[2 dz, and hence that 3t() = 0.\]
\[=\]
\[OCR 2007 9475 Jun07\]
\end{flushleft}
\subsection*{Page 4}
\begin{flushleft}
\bigskip\textbf{Question 8.} (i) Find functions a(r) and b(x) such that u = x and u = e both satisfy the equation\par
\[d2u du\]
\[22 + a(x) da +b(x)u = 0.\]
For these functions a(x) and b(x), write down the general solution of the equation.\par
1 du\par
Show that the substitution y transforms the equation\par
\[= 3u da\]
dy +32+ x 1 (*)\par
\[da +3y + 14 f= 3(1+x)\]
into\par
d²u du 1\par
\[=0\]
\[dn2 + 1+ x dx 1+x\]
\[and hence show that the solution of equation (*) that satisfies y = 0 at x = 0 is given\]
1-e-\par
\[by y =\_3(x+ex)(.\]
\textbf{(ii)} Find the solution of the equation\par
dy x 1\par
+y +\par
dx 1-x 1-x\par
\[that satisfies y = 2 at x = 0.\]
\[OCR 2007 9475 Jun07 [Turn over\]
\end{flushleft}
\subsection*{Page 5}
\begin{flushleft}
\bigskip\textbf{Section B: Mechanics}\par
\bigskip\textbf{Question 9.} Two small beads, A and B, each of mass m, are threaded on a smooth horizontal circular\par
\[hoop of radius a and centre O. The angle 0 is the acute angle determined by 20 = ZAOB.\]
The beads are connected by a light straight spring. The energy stored in the spring is\par
\[mk²a2(0-a)2,\]
\[where k and a are constants satisfying k > 0 and <a <\]
The spring is held in compression with 0 = ẞ and then released. Find the period of oscillations\par
in the two cases that arise according to the value of B and state the value of ẞ for which\par
oscillations do not occur.\par
\bigskip\textbf{Question 10.} A particle is projected from a point on a plane that is inclined at an angle o to the horizontal.\par
The position of the particle at time t after it is projected is (x, y), where (0,0) is the point of\par
projection, measures distance up the line of greatest slope and y measures perpendicular dis-\par
tance from the plane. Initially, the velocity of the particle is given by (t, j) = (V cos 0, V sin0 ),\par
\[where V > 0 and + 0 <ㅠ/2. Write down expressions for x and y.\]
The particle bounces on the plane and returns along the same path to the point of projection.\par
Show that\par
2 tan tan 0 1\par
\[=\]
and that\par
\[V^2 cos2 0\]
\[R=\]
2g sin\par
where R is the range along the plane.\par
Show further that\par
\[2V2\]
\[gR = 3 sin + cosec\]
\[and deduce that the largest possible value of R is V^2/(\sqrt\{3g\}).\]
\[OCR 2007 9475 Jun07\]
\end{flushleft}
\subsection*{Page 6}
\begin{flushleft}
\bigskip\textbf{Question 11.} (i) A wheel consists of a thin light circular rim attached by light spokes of length a to a small\par
hub of mass m. The wheel rolls without slipping on a rough horizontal table directly\par
towards\par
a\par
straight edge of the table. The plane of the wheel is vertical throughout the\par
\[motion. The speed of the wheel is u, where u² < ag.\]
Show that, after the wheel reaches the edge of the table and while it is still in contact\par
with the table, the frictional force on the wheel is zero. Show also that the hub will fall\par
\[vertical distance (ag u^2)/(3g) before the rim loses contact with the table.\]
a -\par
\textbf{(ii)} Two particles, each of mass m/2, are attached to a light circular hoop of radius a, at the\par
ends of a diameter. The hoop rolls without slipping on a rough horizontal table directly\par
towards a straight edge of the table. The plane of the hoop is vertical throughout the\par
motion. When the centre of the hoop is vertically above the edge of the table it has\par
\[speed u, where u^2 < ag, and one particle is vertically above the other.\]
Show that, after the hoop reaches the edge of the table and while it is still in contact\par
with the table, the frictional force on the hoop is non-zero and deduce that the hoop\par
will slip before it loses contact with the table.\par
\[OCR 2007 9475 Jun07 [Turn over\]
\end{flushleft}
\subsection*{Page 7}
\begin{flushleft}
\bigskip\textbf{Section C: Probability and Statistics}\par
\bigskip\textbf{Question 12.} I choose a number from the integers 1, 2,..., (2n - 1) and the outcome is the random\par
\[variable N. Calculate E(N) and E(N^2).\]
I then repeat a certain experiment N times, the outcome of the ith experiment being the\par
\[random variable X; (1 < i N). For each i, the random variable X; has mean u and\]
variance o², and X; is independent of X; for i ≠ j and also independent of N. The random\par
N\par
\[variable Y is defined by Y = Σ X₁. Show that E(Y) = nμ and that Cov (Y, N) = n(n-1)μ.\]
\[Find Var (Y) in terms of n, o^2 and pe.\]
\bigskip\textbf{Question 13.} A frog jumps towards a large pond. Each jump takes the frog either 1 m or 2m nearer to\par
the pond. The probability of a 1 m jump is p and the probability of a 2 m jump is q, where\par
\[p+q= 1, the occurence of long and short jumps being independent.\]
\textbf{(i)} Let Pn(j) be the probability that the frog, starting at a point (n -)m away from\par
the edge of the pond, lands in the pond for the first time on its jth jump. Show that\par
\[p^2(2) p.\]
\[=\]
\textbf{(ii)} Let un be the expected number of jumps, starting at a point (n -) m away from\par
the edge of the pond, required to land in the pond for the first time. Write down the\par
\[value of u^1. By finding first the relevant values of pn(m), calculate u^2 and show that\]
q².\par
\[u^3 =3-2q +\]
\textbf{(iii)} Given that un can be expressed in the form un = A(-q)n-1+B+Cn, where A, B and\par
\[Care constants (independent of n), show that C = (1+q)-1 and find A and B in terms\]
of q. Hence show that, for large n, un and explain carefully why this result is\par
p+2q\par
to be expected.\par
\[OCR 2007 9475 Jun07\]
\end{flushleft}
\subsection*{Page 8}
\begin{flushleft}
14 (i) My favourite dartboard is a disc of unit radius and centre O. I never miss the board,\par
and the probability of my hitting any given area of the dartboard is proportional to the\par
area. Each throw is independent of any other throw. I throw a dart n times (where\par
\[n > 1). Find the expected area of the smallest circle, with centre O, that encloses all\]
the n holes made by my dart.\par
Find also the expected area of the smallest circle, with centre O, that encloses all the\par
(n-1) holes nearest to O.\par
\textbf{(ii)} My other dartboard is a square of side 2 units, with centre Q. I never miss the board,\par
and the probability of my hitting any given area of the dartboard is proportional to the\par
area. Each throw is independent of any other throw. I throw a dart n times (where\par
\[n > 1). Find the expected area of the smallest square, with centre Q, that encloses all\]
the n holes made by my dart.\par
\textbf{(iii)} Determine, without detailed calculations, whether the expected area of the smallest\par
circle, with centre Q, on my square dartboard that encloses all the n holes made by my\par
darts is larger smaller than that for my circular dartboard.\par
or\par
\[OCR 2007 9475 Jun07\]
\end{flushleft}
\end{document}
